\documentclass[11pt]{article}

\usepackage[a4paper,margin=2.6cm]{geometry}
\usepackage{amsmath,amssymb,amsthm}
\usepackage{mathtools}
\usepackage{bm}
\usepackage{array}
\usepackage{booktabs}
\usepackage{xcolor}
\usepackage[colorlinks=true,linkcolor=blue!50!black,citecolor=blue!50!black,urlcolor=blue!50!black]{hyperref}
\usepackage{enumitem}
\usepackage{titlesec}

\titleformat{\section}{\normalfont\Large\bfseries}{\thesection}{0.6em}{}
\titleformat{\subsection}{\normalfont\large\bfseries}{\thesubsection}{0.6em}{}

\newcommand{\dd}{\mathrm{d}}
\newcommand{\pd}{\partial}
\newcommand{\ket}[1]{|#1\rangle}
\newcommand{\bra}[1]{\langle #1|}
\newcommand{\braket}[2]{\langle #1|#2\rangle}
\newcommand{\Op}[1]{\hat{#1}}

\title{\bfseries The Modern State of Global Time in Relativity:\\[2pt] \large A Technical Review}
\author{}
\date{\today}

\begin{document}
\maketitle

\begin{abstract}
\noindent
Relativity replaced the single Newtonian ``now'' with a family of local, frame- and foliation-dependent time functions, and canonical quantum gravity threatens to remove time altogether. Yet physics keeps reconstructing global, or global-\emph{seeming}, time structures at almost every level at which it operates: mathematical relativists prove existence theorems for smooth Cauchy time functions in globally hyperbolic spacetimes; cosmologists use a preferred comoving time singled out by the matter content of the universe; canonical quantum gravity practitioners reconstruct time relationally from matter clocks or from entanglement; algebraic quantum field theorists derive a state-dependent ``thermal time'' from modular flow; the quantum-reference-frame program shows that even proper time and simultaneity become observer- and state-relative operators; parametrized relativistic quantum mechanics (Stueckelberg--Horwitz--Piron theory) reinstates a genuinely universal, though unobservable, evolution parameter; and metrologists have engineered an operational global time scale (TAI/TT) that is accurate to a part in $10^{18}$ and sits consistently on top of curved, dynamical spacetime. This review surveys the current (2024--2026) state of each of these programs, emphasizes the precise sense in which each one does or does not restore a ``global time,'' and offers a comparative synthesis of where the subject presently stands.
\end{abstract}

\tableofcontents

%======================================================================
\section{Introduction}
%======================================================================

Newtonian mechanics assumes a single time parameter $t$, external to and independent of all physical processes, with respect to which every event in the universe can be unambiguously dated and any two events can be unambiguously compared as simultaneous, earlier, or later. Special and general relativity dismantle this picture in stages. Special relativity makes simultaneity frame-dependent: two events simultaneous for one inertial observer are generally not simultaneous for another. General relativity goes further, making the very existence of \emph{any} global time-like structure a nontrivial, spacetime-dependent, and sometimes false, mathematical proposition. Canonical approaches to quantum gravity go further still: in the Wheeler--DeWitt formulation the total Hamiltonian is a constraint that annihilates physical states, so that, formally, nothing evolves in the bulk of the theory at all. This sequence of moves, from a single external clock to no clock whatsoever, is often summarized as the ``problem of time.''

At the same time, essentially every branch of physics that touches relativity has needed to reconstruct \emph{some} global or quasi-global temporal structure to make contact with observation, prediction, and everyday practice: a preferred slicing for numerical relativity and cosmology, a well-defined ``clock reading'' for canonical quantum gravity, a single planetary time scale for satellite navigation and geodesy. The interesting modern content is therefore not a single yes/no verdict on whether global time exists, but a differentiated map of \emph{which} notion of global time survives, in \emph{which} regime, under \emph{which} conditions, and at \emph{what} price.

This review builds that map. Section~\ref{sec:taxonomy} lays out a taxonomy of distinct things that can be meant by ``global time.'' Section~\ref{sec:sr} recalls why special relativity already forecloses the Newtonian notion. Section~\ref{sec:gr} treats the mathematically precise question of when a classical, generally covariant spacetime admits a global time function, via the Geroch and Bernal--S\'anchez splitting theorems, the ADM 3+1 formalism, and cosmological time. Section~\ref{sec:pot} reviews the problem of time in canonical quantum gravity and its Kuch\v{a}\v{r}--Isham classification, updated with 2024--2026 technical work. Section~\ref{sec:relational} covers relational and emergent time, centered on the Page--Wootters mechanism and its recent generalizations, tests, and equivalence results. Section~\ref{sec:thermal} covers the Connes--Rovelli thermal time hypothesis. Section~\ref{sec:qrf} covers the quantum-reference-frame program, in which proper time and simultaneity themselves become relative, state-dependent quantum observables, with genuine experimental proposals now on the table. Section~\ref{sec:shp} treats Stueckelberg--Horwitz--Piron (SHP) theory, a parametrized relativistic quantum theory in which a strictly external, universal, monotonically increasing invariant parameter $\tau$ is reinstated as the fundamental evolution time, at the cost of demoting ordinary coordinate time $x^0$ to a dynamical variable. Section~\ref{sec:metrology} turns to the practical, operational side: how a single global time scale (TAI/TT/TCG) is actually constructed and disseminated on top of a fully relativistic Earth, and where that program stands as the SI second heads toward redefinition. Section~\ref{sec:synthesis} synthesizes the results in a comparative table, and Section~\ref{sec:outlook} closes with open problems.

%======================================================================
\section{A Taxonomy of ``Global Time''}
\label{sec:taxonomy}
%======================================================================

Before assessing claims about global time it is worth separating several logically distinct questions that the phrase can refer to:

\begin{enumerate}[label=(\alph*)]
\item \textbf{Absolute simultaneity.} Is there a unique, frame-independent relation of ``at the same time'' between spatially separated events? (Newtonian: yes. Special relativity: no.)
\item \textbf{A global time function.} Does a given spacetime $(M,g)$ admit a smooth scalar field $t:M\to\mathbb{R}$ with everywhere past- or future-pointing timelike gradient, whose level sets foliate $M$ into spacelike slices? This is a well-posed differential-geometric question with a definite theorem-based answer for each spacetime (Section~\ref{sec:gr}).
\item \textbf{A preferred (canonical) global time.} Given that a global time function exists, is one particular choice of slicing singled out by the physics (matter content, symmetry, or dynamics), rather than being a free gauge choice among infinitely many?
\item \textbf{A time variable in the quantum theory.} Once the classical theory is quantized in a generally covariant way, is there any variable, external or internal, with respect to which the physical states evolve, given that the total (super-)Hamiltonian annihilates them?
\item \textbf{A universal evolution parameter.} Is there a single, Lorentz-invariant, monotonically increasing parameter with respect to which \emph{all} events/particles in the universe evolve, playing the role Newton's $t$ played, even though it need not coincide with any observable coordinate or proper time?
\item \textbf{An operational global time standard.} Is there an internationally agreed, physically realizable time scale that laboratories, satellites, and observatories around the Earth can synchronize to, accounting for all relativistic effects?
\end{enumerate}

As will become clear, the modern literature answers (a) negatively, (b) conditionally (``yes, if and only if the spacetime is globally hyperbolic''), (c) contextually (cosmology says approximately yes; generic solutions say no), (d) is the crux of the problem of time and has several competing, partial resolutions, (e) is answered affirmatively by construction in SHP-type parametrized theories, and (f) is answered pragmatically and very successfully by relativistic metrology. Keeping these six senses distinct is essential, because a great deal of both technical and popular confusion arises from equivocating between them.

%======================================================================
\section{The Relativity of Simultaneity}
\label{sec:sr}
%======================================================================

\subsection{Einstein synchronization and its frame-dependence}

In special relativity, an inertial observer assigns simultaneity to distant events by Einstein's radar/light-signal convention: a light signal is sent at (local) time $t_1$, reflected at the distant event, and received back at $t_2$; the event of reflection is assigned time $t_1+\tfrac12(t_2-t_1)$. Under a Lorentz boost with velocity $v$ along $x$, coordinates transform as
\begin{equation}
t' = \gamma\left(t - \frac{v x}{c^2}\right), \qquad x' = \gamma\left(x - v t\right), \qquad \gamma = \left(1-\frac{v^2}{c^2}\right)^{-1/2}.
\end{equation}
Two events with $\Delta t = 0$ but $\Delta x \neq 0$ in one frame therefore satisfy $\Delta t' = -\gamma v \Delta x/c^2 \neq 0$ in a boosted frame: simultaneity for spatially separated events is not preserved under a change of inertial frame. There is no invariant meaning to ``now, over there.'' The invariant structure of Minkowski spacetime is instead the light-cone (causal) structure, encoded in the interval
\begin{equation}
\dd s^2 = -c^2\dd t^2 + \dd x^2+\dd y^2+\dd z^2,
\end{equation}
together with the causal order it induces; simultaneity surfaces are merely one family of spacelike slices among infinitely many, each tied to a choice of inertial frame.

\subsection{The conventionality debate and Malament's theorem}

Reichenbach observed that Einstein's ``half-way'' convention is not forced by the causal structure alone: replacing the factor $\tfrac12$ by an arbitrary $\varepsilon\in(0,1)$ in $t_1+\varepsilon(t_2-t_1)$ yields a different, still logically consistent, simultaneity relation, and asked whether the standard choice $\varepsilon=1/2$ is a substantive physical fact or a mere convention. Malament's 1977 theorem sharpened this debate considerably: he showed that, within a single inertial frame of Minkowski spacetime, standard ($\varepsilon=1/2$) simultaneity is in fact the \emph{unique} nontrivial equivalence relation on events definable purely from the causal (light-cone) structure that reduces, for a given observer, to the ordinary co-located ``instant'' relation. This is widely read as showing that once one fixes an inertial frame, simultaneity within that frame is not an arbitrary convention but is fixed by the causal structure together with the choice of frame; the residual conventionality debate in the philosophy-of-physics literature concerns exactly how much this settles, since the theorem still presupposes the selection of a rest frame, and different responses to it (Janis, Norton, and others) continue to be discussed. What survives uncontroversially, in any case, is the frame-\emph{relativity} of simultaneity itself, which is precisely what removes Newtonian global time from special relativity.

%======================================================================
\section{Global Time Functions in Curved Spacetime}
\label{sec:gr}
%======================================================================

General relativity replaces the fixed Minkowski background with a dynamical Lorentzian manifold $(M,g)$, and the question ``does a global time exist'' becomes the mathematical question of whether $M$ admits a global \emph{time function}.

\subsection{Global hyperbolicity and the splitting theorems}
\label{ssec:gh}

A spacetime is \emph{globally hyperbolic} if it is causal (no closed causal curves) and the causal past-and-future intersection $J^+(p)\cap J^-(q)$ of any two points is compact. Global hyperbolicity is the condition under which the Einstein equations admit a well-posed initial value problem: initial data on a single \emph{Cauchy surface} $\Sigma$ (an achronal set met exactly once by every inextendible causal curve) determine the entire spacetime.

Geroch (1970) proved that a spacetime is globally hyperbolic if and only if it admits a Cauchy surface, and that in this case $M$ is homeomorphic to $\mathbb{R}\times\Sigma$, with a continuous \emph{time function} whose level sets are (topological) Cauchy surfaces. This settled the topological question but left the smooth (differentiable) version open for over three decades. Bernal and S\'anchez closed this gap: they showed that every globally hyperbolic spacetime admits a \emph{smooth} spacelike Cauchy hypersurface, hence a smooth global diffeomorphism $M\cong \mathbb{R}\times\Sigma$, and moreover that it admits a smooth \emph{Cauchy temporal function} $\tau$ (a smooth function with past-directed timelike gradient everywhere) whose level sets foliate $M$ by smooth Cauchy surfaces, so that the metric can be written in the split form
\begin{equation}
g = -\beta\, \dd\tau^2 + \bar g_\tau, \qquad \bar g_\tau \ \text{a Riemannian metric on each slice } \Sigma_\tau,
\label{eq:splitting}
\end{equation}
with $\beta>0$ a smooth function. This is the modern, fully rigorous statement of ``a global time function exists'' in classical general relativity: it is available whenever, and only whenever, the spacetime is globally hyperbolic, which is exactly the class of spacetimes for which Einstein's equations are predictively well-posed in the first place. Related work has since characterized global hyperbolicity intrinsically via the completeness of a \emph{null distance} built from a temporal function, and via Lorentzian analogues of Cauchy--Riemannian length-space geometry, extending the splitting theorem to low-regularity and even non-smooth (Lorentzian length space) settings.

\subsection{Many-fingered time: the ADM decomposition}
\label{ssec:adm}

Once a spacelike foliation $\{\Sigma_t\}$ is chosen, the Arnowitt--Deser--Misner (ADM) formalism writes the metric as
\begin{equation}
\dd s^2 = -N^2\dd t^2 + h_{ij}\left(\dd x^i + N^i \dd t\right)\left(\dd x^j + N^j \dd t\right),
\label{eq:adm}
\end{equation}
where $h_{ij}$ is the induced spatial metric on $\Sigma_t$, $N$ is the \emph{lapse} function (the proper time elapsed between neighboring slices, per unit coordinate time, for an observer moving orthogonally to the slices), and $N^i$ is the \emph{shift} vector (how spatial coordinates are dragged between slices). Crucially, $N$ and $N^i$ are not fixed by the dynamics; they are freely specifiable gauge functions, reflecting the arbitrariness of \emph{which} foliation and \emph{which} spatial coordinates one has chosen. Wheeler termed this freedom ``many-fingered time'': at every point of space, the rate and even direction of the local time flow relative to the global label $t$ can be independently dialed, subject only to $\Sigma_t$ remaining spacelike. A global time function of the Bernal--S\'anchez type exists, in this language, precisely when a lapse and shift can be chosen so that the foliation covers all of $M$ without the slices ever failing to be Cauchy surfaces; but there is an infinite-dimensional family of equally valid choices, so ``a'' global time function is never ``the'' global time function.

\subsection{Failures of global time}

Not every physically interesting spacetime is globally hyperbolic. The G\"odel universe (a rotating dust cosmology) contains closed timelike curves and admits no time function of any kind; the maximally extended Kerr and Reissner--Nordstr\"om black-hole interiors contain Cauchy horizons beyond which predictability, and with it any global time function, breaks down; and spacetimes with ``naked'' causality violations or with incomplete, non-compactly-generated horizons can similarly fail global hyperbolicity. In these cases the question of global time does not have a qualified answer -- it has the answer \emph{no}, full stop, illustrating that even at the purely classical level a universal time is a contingent, not guaranteed, feature of relativistic spacetimes.

\subsection{Cosmological time}

Physically realistic cosmological spacetimes, by contrast, are exceptionally well-behaved. The homogeneous, isotropic Friedmann--Lema\^itre--Robertson--Walker (FLRW) metric,
\begin{equation}
\dd s^2 = -c^2\dd t^2 + a(t)^2\left[\frac{\dd r^2}{1-kr^2}+r^2\dd\Omega^2\right],
\end{equation}
comes with an essentially canonical global time function: the cosmic (comoving) time $t$, defined as proper time along the worldlines of observers comoving with the average matter distribution, whose constant-$t$ slices are homogeneous and isotropic. This is not merely a mathematically convenient gauge choice; it is singled out physically by the symmetry of the matter content (equivalently, by the local rest frame of the cosmic microwave background), and it agrees with the cosmological time function constructed intrinsically from the causal structure of spacetimes with (past-)crushing singularities. It is in this sense the strongest real-world example of case (c) in the taxonomy of Section~\ref{sec:taxonomy}: a global time function that is preferred, not merely available.

%======================================================================
\section{The Problem of Time in Canonical Quantum Gravity}
\label{sec:pot}
%======================================================================

\subsection{The Hamiltonian constraint and the frozen formalism}

Applying the ADM split to canonical general relativity produces a Hamiltonian that is a sum of constraints,
\begin{equation}
H = \int_\Sigma \dd^3x \, \left( N \mathcal{H}_\perp + N^i \mathcal{H}_i \right),
\end{equation}
with the Hamiltonian (scalar) constraint $\mathcal{H}_\perp\approx 0$ and the momentum (diffeomorphism) constraints $\mathcal{H}_i\approx 0$ holding as identities on the physical phase space, a direct consequence of the reparametrization/diffeomorphism invariance of general relativity: the lapse and shift are Lagrange multipliers, not dynamical fields, so the numerical value of $H$ on-shell is zero. Dirac quantization promotes $\mathcal{H}_\perp$ to an operator and demands that physical states be annihilated by it,
\begin{equation}
\hat{\mathcal{H}}_\perp\,\Psi[h_{ij}] = 0,
\end{equation}
the Wheeler--DeWitt equation. Because the total Hamiltonian vanishes on physical states, there is no Schr\"odinger-type generator left over to produce time evolution: $\Psi$ satisfies a single, time-independent equation on the space of 3-geometries (superspace), and is sometimes described as the (static) ``wavefunction of the universe.'' This is the technical core of the problem of time: general covariance, taken seriously into the quantum theory, removes the external clock with respect to which Schr\"odinger evolution is normally defined.

\subsection{The Kuch\v{a}\v{r}--Isham classification}

The influential early-1990s reviews by Kuch\v a\v r and by Isham organized responses into three broad families, and this taxonomy remains the standard reference point:
\begin{itemize}
\item \textbf{Time before quantization.} Identify a physical degree of freedom already at the classical level -- a matter field, a scalar clock, or an intrinsic geometric quantity such as the local volume or extrinsic curvature -- to serve as an internal clock, then \emph{deparametrize} with respect to it before quantizing, producing an ordinary (non-relativistic-looking) Schr\"odinger equation for the remaining degrees of freedom. This trades general covariance (clock choice becomes a gauge-fixing) for a genuine unitary evolution.
\item \textbf{Time after quantization.} Quantize the full constrained system first, and only afterward try to recover an approximate, emergent semiclassical time, typically via a Born--Oppenheimer-type separation between ``heavy'' gravitational and ``light'' matter degrees of freedom, or via a WKB expansion of the Wheeler--DeWitt wavefunctional in $\hbar$.
\item \textbf{Timeless strategies.} Accept that there is no fundamental time and reinterpret physics accordingly: the naive Schr\"odinger interpretation (treat $|\Psi[h_{ij}]|^2$ directly as a static probability), the conditional-probability interpretation (ask for correlations between subsystems rather than time evolution of a single system -- the direction taken by the Page--Wootters mechanism, Section~\ref{sec:relational}), records theories (treat the present state as containing ``records'' of a history without there being any actual flow), and Rovelli's programme of \emph{partial and complete observables} and \emph{evolving constants of the motion}, in which one computes the correlation between a chosen internal clock variable and any other observable, without ever privileging one variable as ``the'' time.
\end{itemize}

\subsection{Recent technical work (2024--2026)}

The problem of time remains an active research area rather than a closed textbook topic. Representative recent contributions include: geometric constructions of a ``clock field'' tied to the extrinsic curvature of the foliation and the normal congruence, designed to identify precisely the geometric conditions under which a relational description of time is or is not valid, without adding new degrees of freedom to the Einstein--Hilbert action; reassessments arguing that the standard resolutions of the problem of time rest on questionable foundational assumptions shared by geometrodynamics, loop quantum gravity, and quantum cosmology alike; and essay-level treatments arguing that the vanishing of the Hamiltonian should be read not as a defect but as a structural feature (``projectivity'') of the physical Hilbert space in quantum gravity, connected to the existence of an entire vacuum module rather than a single vacuum state. This diversity of live positions -- some seeking to repair internal/matter time constructions, some doubting that any repair is possible, some reinterpreting the vanishing Hamiltonian as a feature -- is itself the most accurate one-line summary of where the field stands: there is technical consensus on the \emph{mathematical origin} of the problem (the linear combination of first-class constraints forming $H$), and continued, active disagreement about its \emph{physical resolution}.

%======================================================================
\section{Relational and Emergent Time}
\label{sec:relational}
%======================================================================

\subsection{Deparametrization and the Page--Wootters mechanism}

The most influential ``timeless'' construction is due to Page and Wootters (1983): partition the total Hilbert space into a clock subsystem $C$ and the rest of the universe $R$, and impose the Wheeler--DeWitt-type constraint on the global state,
\begin{equation}
\left(\Op{H}_C\otimes \mathbb{1}_R + \mathbb{1}_C\otimes \Op{H}_R\right)\ket{\Psi} = 0.
\end{equation}
Although $\ket{\Psi}$ is, by construction, a stationary (``frozen'') state of the total constraint, conditioning on the clock reading a particular value $\tau$,
\begin{equation}
\ket{\psi_R(\tau)} \ \propto\ \left(\bra{\tau}_C\otimes\mathbb{1}_R\right)\ket{\Psi},
\end{equation}
yields, under mild conditions, an effective Schr\"odinger evolution $i\hbar\,\pd_\tau\ket{\psi_R(\tau)} = \Op{H}_R\ket{\psi_R(\tau)}$ for the rest of the universe relative to the clock. Time, on this reading, is not a fundamental background parameter but a relational property that emerges from correlations between subsystems of an otherwise static entangled state; an external observer of the whole (clock $+$ rest) sees no evolution at all, while an observer with access only to the clock-conditioned branch sees ordinary quantum dynamics. Early technical objections -- concerning multi-time correlations and the physical meaning of the clock ``reading'' -- have substantially been addressed by more careful reformulations, and a body of work now known informally as the ``trinity of relational dynamics'' has established that the Page--Wootters formalism, Dirac (relational) quantization, and classical deparametrize-then-quantize approaches are, under appropriate conditions, equivalent descriptions of the same physics, including in settings with periodic or otherwise non-monotonic clocks.

\subsection{Extensions, group field theory, and non-perturbative applications}

The formalism continues to be extended in directions directly relevant to full, non-perturbative quantum gravity. It has recently been applied to group field theory, a framework in which spacetime itself, rather than being fixed in advance, is meant to emerge from an underlying combinatorial/field-theoretic structure and in which there is consequently no background time of any kind by construction; there the Page--Wootters approach has been shown to connect to \emph{multi-fingered time} (Section~\ref{ssec:adm}) generalized to many independent gauge symmetries. Related work addresses the ``time-of-arrival'' problem -- when, relationally, does a given event occur -- within the Page--Wootters formalism, and clarifies the regularization needed to obtain well-defined conditional probabilities for that question.

\subsection{Experimental and thermodynamic probes}

What was, until recently, an entirely formal debate now has an emerging experimental dimension. Recent work on quantum clock thermodynamics has shown that extracting a usable, classical-looking sequence of clock ``ticks'' from an underlying quantum clock carries an intrinsic thermodynamic and informational cost -- in one 2025 study using single electrons hopping between quantum dots as an explicit clock, the energetic cost of \emph{reading} the clock was found to be many orders of magnitude larger than the cost of simply running it, directly probing the entropic price of converting quantum, relational time into the familiar classical ticking we experience. In a related direction, toy-model studies building on the original Page--Wootters construction have shown how entanglement between subsystems of an initially timeless, spaceless quantum state can generate not only an effective flow of time for a conditioned observer but also an effective sense of spatial order, suggesting that relational/entanglement-based emergence may extend beyond time to spatial structure itself.

%======================================================================
\section{Thermal Time}
\label{sec:thermal}
%======================================================================

A structurally different response to the problem of time, due to Connes and Rovelli, notes that a generic statistical (mixed) state $\rho$ on a classical phase space, or a faithful normal state $\omega$ on a von Neumann algebra of observables in the quantum case, comes equipped with a canonically associated one-parameter flow even in a generally covariant theory with no a priori time variable. Classically, given $\rho$, define the \emph{thermal Hamiltonian} $K_\rho \equiv -\ln\rho$ and let \emph{thermal time} be the parameter of the flow generated by $K_\rho$ on phase space. Quantum mechanically, the analogous flow is the \emph{modular flow} $\sigma^\omega_t$ canonically attached to $\omega$ by Tomita--Takesaki theory, which is automatically a one-parameter group of automorphisms of the observable algebra satisfying the Kubo--Martin--Schwinger (KMS) equilibrium condition with respect to itself at inverse temperature $\beta=1$. The thermal time hypothesis is the proposal that this modular flow, rather than being merely a mathematical curiosity of equilibrium statistical mechanics, \emph{is} the physical flow of time perceived by a system when it is found in the state $\omega$: time, on this view, is not put in by hand but is thermodynamically induced by whatever statistical/coarse-grained state the system happens to occupy, and different states of the same generally covariant system can induce different thermal times.

The hypothesis passes a striking consistency check: in an FLRW universe, the thermal time associated with the thermal state of the cosmic microwave background radiation has been shown to coincide with the ordinary Robertson--Walker cosmic time of Section~\ref{sec:gr}, recovering the standard cosmological clock from a purely thermodynamic construction rather than assuming it. More recent mathematical work has clarified the operational status of thermal time by constructing an explicit positive-operator-valued measure (POVM), i.e., an ``unsharp'' quantum time observable, covariant with the modular flow for concrete systems such as the quantum harmonic oscillator, putting the hypothesis on a more standard measurement-theoretic footing. At the same time, the thermal time hypothesis has attracted sustained philosophical and technical criticism: objections include the unclear relationship between thermal time and ordinary proper time away from the special cosmological example, doubts about whether the construction is truly background-independent and gauge-invariant once matter and gravitational degrees of freedom are both included, and the question of what, if anything, thermal time means for systems that are not in (or are far from) statistical equilibrium.

%======================================================================
\section{Quantum Reference Frames: Simultaneity and Proper Time as Quantum Observables}
\label{sec:qrf}
%======================================================================

A third, largely independent research program -- quantum reference frames (QRFs) -- asks what happens to relativistic notions of time when the reference frame itself, or the clock used to define proper time, is allowed to be in a quantum superposition, rather than simply choosing between one classical frame and another.

\subsection{Superpositions of proper time}

If an internal clock degree of freedom is attached to a particle moving through curved spacetime along a trajectory that is itself in a quantum superposition of two paths with different proper-time elapse $\tau_1\ne\tau_2$ (for instance, two arms of a matter-wave interferometer at different heights in a gravitational field), the reduced state of the clock becomes entangled with the spatial degree of freedom, and the resulting loss of interferometric visibility becomes a direct, in-principle measurable signature of a genuine \emph{superposition of proper times} rather than of a classical mixture of them. This produces a proper-time--mass/energy uncertainty relation completely analogous to the ordinary energy--time uncertainty relation, but derived, this time, from the relativistic structure of proper time itself treated as a covariant quantum observable (formally a POVM rather than a self-adjoint operator, since proper time, like ordinary time, is not conjugate to a bounded-below Hamiltonian in the usual Stone--von Neumann sense). Concrete experimental proposals now target this effect directly: trapped-atom interferometers held in a superposition of heights for on the order of a minute, combined with auxiliary clock transitions, have been proposed as a near-term platform to witness a superposition of relativistic proper times via both visibility modulation and a shift in resonant frequency; a closely related 2026 proposal targets the same physics using optical ion clocks.

\subsection{Indefinite causal order, quantum coordinates, and the relativity of event localization}

A broader program, associated especially with Brukner and collaborators, develops a full covariant formalism for transforming between reference frames that are themselves quantum systems (e.g., massive bodies in a superposition of trajectories), generalizing both special-relativistic frame transformations and the equivalence principle to this setting. A striking upshot is that even the \emph{localization of an event} -- where and when something happens -- becomes frame-relative in a way that goes beyond ordinary relativity: two quantum reference frames can disagree not merely on the time-coordinate of an event (as in special relativity) but on which events are localized at all, and recent work formalizes this using ``quantum coordinates'' and an explicit quantum analogue of Einstein's hole argument. Related developments extend these transformations to indefinite causal order, in which the very question of which of two events happens first can become quantum-superposed rather than merely frame-dependent, and to settings with an indefinite spacetime metric, where even the light-cone structure that ordinarily defines simultaneity and causal order is itself put in superposition. Taken together, this body of work shows that the relativity of simultaneity established by special relativity (Section~\ref{sec:sr}) is not the end of the story: at the operational, information-theoretic level accessible to a genuinely quantum observer, simultaneity, proper time, and even event identity all become state-dependent quantum variables rather than classical, frame-relative but otherwise well-defined, quantities.

%======================================================================
\section{A Genuinely Universal Time: Stueckelberg--Horwitz--Piron Theory}
\label{sec:shp}
%======================================================================

Every construction surveyed so far either denies a fundamental global time outright (canonical quantum gravity, thermal time as state-dependent, quantum reference frames) or recovers one only conditionally, locally, or approximately (Cauchy time functions, cosmic time, relational/matter time). Stueckelberg--Horwitz--Piron (SHP) theory takes the opposite strategy: it reinstates, by explicit postulate, a single external, Lorentz-invariant, monotonically increasing evolution parameter $\tau$ that plays exactly the role of Newton's absolute time -- at the price of demoting the ordinary laboratory time coordinate $x^0$ from an external parameter to a fully dynamical variable on equal footing with the three spatial coordinates.

\subsection{Motivation: pair creation and the failure of $x^0$ and proper time as evolution parameters}

The construction originates with Fock's 1937 observation that proper time can serve as an evolution parameter for a manifestly covariant, Newton-like force law, and with Stueckelberg's 1941 reinterpretation of antiparticles as ordinary particles whose worldlines locally turn around and propagate backward in \emph{coordinate} time $x^0$, later reformulated by Feynman. Stueckelberg noted that neither $x^0$ nor the ordinary proper time $s$, with $\dd s^2 = -\eta_{\mu\nu}\dd x^\mu\dd x^\nu$, can consistently serve as \emph{the} evolution parameter once such pair creation/annihilation processes are allowed, since $\dd s^2/\dd\tau^2$ is not sign-definite along a worldline that turns around in time and $x^0$ itself is not monotonic along such a trajectory. He therefore introduced an external, strictly monotonic chronological parameter $\tau$, distinct from both, with respect to which the four spacetime coordinates $x^\mu(\tau)$ (including $x^0(\tau)$) evolve as a single unconstrained dynamical trajectory. Horwitz and Piron generalized this in 1973 into a fully relativistic many-body theory, in which $\tau$ is promoted to a genuinely universal invariant time shared by every particle in the system -- structurally the direct relativistic analogue of the single Newtonian $t$, but now playing the role of an unconstrained, off-shell evolution parameter for events in an 8-dimensional single-particle phase space $(x^\mu,p_\mu)$, rather than of a coordinate on spacetime itself.

\subsection{Formalism: off-shell dynamics and $\tau$-dependent gauge fields}

The free single-event Hamiltonian is taken to be Lorentz-invariant and quadratic in momentum,
\begin{equation}
K_0 = \frac{\eta^{\mu\nu}p_\mu p_\nu}{2M},
\end{equation}
generating $\tau$-evolution via Hamilton's equations $\dd x^\mu/\dd\tau = \pd K/\pd p_\mu$, $\dd p_\mu/\dd\tau = -\pd K/\pd x^\mu$, with $M$ a fixed parameter of dimension mass that need not coincide with the particle's observed (on-shell) mass; the mass-shell condition $p^\mu p_\mu = -M^2c^2$ of ordinary relativistic mechanics is recovered only asymptotically, in a $\tau$-independent equilibrium limit, rather than being imposed as a constraint from the outset. Interactions are introduced through a set of five $\tau$-dependent gauge potentials $a^\mu(x,\tau)$ (four spacetime components plus one component conjugate to $\tau$ itself), generalizing electrodynamics to a manifestly $\tau$-local $U(1)$ gauge theory (``pre-Maxwell'' theory), with field strengths $f^{\mu\nu}=\pd^\mu a^\nu - \pd^\nu a^\mu$ that reduce to the ordinary Maxwell field only once the system relaxes to a $\tau$-independent equilibrium. This framework has been used, among other things, to give a natural covariant account of electron interference-in-time experiments and to construct a relativistic Boltzmann equation with a genuine, single, universal evolution parameter and an associated entropy flow, avoiding certain irreversibility paradoxes that arise when $x^0$ itself is used as the statistical evolution parameter.

At the level of quantum theory, the state of a system is a wavefunction $\psi(x,\tau)$ obeying a Stueckelberg--Schr\"odinger-type equation $i\hbar\,\pd_\tau\psi = \Op{K}\psi$ directly analogous to ordinary non-relativistic quantum mechanics, but with $x^0$ now an operator on the same footing as $x^1,x^2,x^3$, rather than an external label -- a structural move that several authors have noted is formally close to, and historically anticipates, aspects of the relational/deparametrized constructions of Section~\ref{sec:relational}, although SHP theory retains $\tau$ itself as a strictly external, non-dynamical, unobservable label, whereas fully relational programs insist that even $\tau$-like labels must ultimately be internal.

\subsection{Extension to curved spacetime and open issues}

SHP theory has been extended into a full general-relativistic setting (``SHPGR''), in which the canonical Poisson-bracket structure of the flat-space theory is shown to remain invariant under general coordinate transformations on a curved background, providing a basis for a corresponding canonical quantum theory, and in which the scalar potential of the flat-space theory reappears as an induced spacetime mass distribution. A complementary ``4+1'' formalism generalizes the ADM decomposition itself (Section~\ref{ssec:adm}) to an evolving five-dimensional local metric $g_{\mu\nu}(x,\tau)$, in which the four-dimensional spacetime metric is dynamical not only in the usual GR sense but also \emph{evolves with $\tau$}, with the standard 4D Einstein equations recovered only in a $\tau$-independent limit; this construction deliberately breaks the naive five-dimensional symmetry of the formalism to avoid collapsing into ordinary (Kaluza--Klein-type) 5D gravity. Open technical issues include the causal structure of the associated five-dimensional Green's functions, which develop singularities tied to $\tau$-dependence that must be regularized (e.g., by adding higher-order kinetic terms for the gauge fields) while preserving gauge and Lorentz invariance; the resulting regularized electrodynamics is super-renormalizable, which is itself of some field-theoretic interest independent of the time-foundations motivation. In this framework, then, questions (b)--(d) of the taxonomy in Section~\ref{sec:taxonomy} are answered essentially as in standard GR (a global \emph{coordinate}-time function may or may not exist, and is never uniquely preferred), while question (e) -- is there a universal evolution parameter -- is answered affirmatively by explicit construction, precisely by refusing to identify that parameter with any observable spacetime coordinate.

%======================================================================
\section{The Practical Realization of Global Time: Relativistic Metrology}
\label{sec:metrology}
%======================================================================

Independently of how the foundational debates above are eventually resolved, the world runs on an operational global time standard that is fully relativistic by construction, and the precision of that standard has been advancing rapidly.

\subsection{TAI, Terrestrial Time, and coordinate times}

International Atomic Time (TAI) is a weighted statistical average over several hundred atomic clocks in national metrology laboratories worldwide, steered to the SI second as realized on the rotating geoid. Terrestrial Time (TT), used as the independent argument for geocentric ephemerides, is defined with a fixed offset from TAI (chosen, historically, to match the old ephemeris time scale at the point of transition) and is, by construction, the proper time that would be kept by an ideal clock at rest on the geoid, correcting for the gravitational and rotational time dilation appropriate to that surface. Two further, purely coordinate (as opposed to proper) time scales complete the standard relativistic framework recommended by the International Astronomical Union: Geocentric Coordinate Time (TCG), the coordinate time of a reference frame centered on the Earth but not corotating with, or gravitationally bound to, its surface, and Barycentric Coordinate Time (TCB), the analogous coordinate time centered on the solar-system barycenter. TCG and TT differ by a small, precisely defined constant rate, reflecting the gravitational and rotational redshift between the geoid and a notional observer at Earth's center; TCB and TT differ by both a constant rate and a periodic (annual) term arising from the Earth's motion in the Sun's gravitational potential. This four-tiered structure (TAI, TT, TCG, TCB) is, in effect, an explicit, standards-body-ratified answer to question (f) of the taxonomy: a single, global, operationally realizable time convention, built by choosing one physically motivated foliation (the geoid, or the solar-system barycentric frame) among the many allowed by general relativity, and then correcting every local proper-time reading to that convention via well-tested relativistic formulas -- precisely the ``many-fingered time'' freedom of Section~\ref{ssec:adm}, resolved by explicit international convention rather than by any claim of physical uniqueness.

\subsection{Chronometric geodesy and optical clock networks}

The precision of optical atomic clocks (based on optical, rather than microwave, transitions in ions or neutral atoms held in optical lattices) has now reached fractional frequency uncertainties below one part in $10^{18}$, sufficient to resolve the general-relativistic gravitational redshift between two clocks separated by only centimeters of height. This has founded the field of \emph{chronometric} (or relativistic) geodesy, in which differences in clock rate are used directly to map the Earth's gravitational (geo)potential, complementing traditional gravimetric and satellite geodesy. Coordinated international comparisons have recently linked ten next-generation optical clocks across six countries via fiber and satellite links, and dedicated space missions (e.g., an atomic-clock ensemble on the International Space Station) aim to push relativistic-redshift tests, and the associated geopotential mapping, substantially further. These networks are simultaneously scientific instruments, tests of general relativity, and the technical infrastructure through which a future all-optical global time scale would be realized and disseminated.

\subsection{Toward redefining the second, and retiring the leap second}

The SI second is still officially defined via the microwave hyperfine transition of caesium-133, but international metrology bodies (coordinated through the Consultative Committee for Time and Frequency and the General Conference on Weights and Measures, CGPM) have been developing a roadmap toward redefining the second in terms of an optical transition once several independent optical clocks worldwide demonstrate the required accuracy and mutual agreement, with a decision path currently aimed at the early 2030s. In parallel, the 27th CGPM (2022) resolved to eliminate the leap second from UTC by, at the latest, 2035, removing an irregular, unpredictable correction that had itself been a direct, practical consequence of the mismatch between a uniform relativistic time scale (TAI) and the non-uniform rotation of the Earth. Both developments illustrate the same underlying point from a metrological angle: a single, global, practically usable time scale is entirely achievable in a relativistic universe, but only as an explicit engineering convention, continuously corrected for well-understood relativistic effects, never as a rediscovery of Newtonian absolute time.

%======================================================================
\section{Synthesis}
\label{sec:synthesis}
%======================================================================

Table~\ref{tab:synthesis} collects the verdicts reached above for each sense of ``global time'' distinguished in Section~\ref{sec:taxonomy}.

\begin{table}[htbp]
\centering
\small
\renewcommand{\arraystretch}{1.35}
\begin{tabular}{@{}p{3.6cm}p{5.6cm}p{5.6cm}@{}}
\toprule
\textbf{Notion of global time} & \textbf{Status} & \textbf{Where it lives} \\
\midrule
Newtonian absolute simultaneity & Refuted & Pre-relativistic mechanics \\
Global time function on a fixed spacetime & Exists iff globally hyperbolic (Geroch; Bernal--S\'anchez); non-unique (``many-fingered'') & Classical GR (Sec.~\ref{sec:gr}) \\
Preferred cosmological time & Physically singled out by matter content / CMB rest frame & FLRW cosmology (Sec.~\ref{sec:gr}) \\
Time in canonical quantum gravity & Absent in the fundamental (Wheeler--DeWitt) formalism -- the ``problem of time'' & Quantum geometrodynamics, loop quantum gravity (Sec.~\ref{sec:pot}) \\
Relational / matter time & Reconstructed via a chosen internal clock subsystem; not unique, approximate away from idealized clocks & Page--Wootters, deparametrization (Sec.~\ref{sec:relational}) \\
Thermal (modular) time & State-dependent; coincides with cosmic time for the CMB thermal state; contested away from equilibrium & Algebraic QFT, Connes--Rovelli hypothesis (Sec.~\ref{sec:thermal}) \\
Simultaneity / proper time under quantum reference frames & Becomes a state-dependent quantum observable; can be put in superposition; event localization itself is frame-relative & Quantum-reference-frame program (Sec.~\ref{sec:qrf}) \\
Universal invariant evolution parameter $\tau$ & Postulated by construction; external and non-observable; $x^0$ becomes dynamical & Stueckelberg--Horwitz--Piron theory (Sec.~\ref{sec:shp}) \\
Operational global time standard & Achieved to $\sim 10^{-18}$ precision by explicit international convention (TAI/TT/TCG/TCB), continuously relativistically corrected & Metrology, geodesy, satellite navigation (Sec.~\ref{sec:metrology}) \\
\bottomrule
\end{tabular}
\caption{A comparative summary of the distinct notions of ``global time'' surveyed in this review, and their current status.}
\label{tab:synthesis}
\end{table}

The throughline is that general relativity permits a global time function only conditionally (global hyperbolicity) and never uniquely (many-fingered time); canonical quantization removes even that conditional structure; and every one of the major modern responses -- relational/matter time, thermal time, quantum-reference-frame constructions, and SHP-type universal parameters -- recovers an effective or reinterpreted notion of time only by adding an extra structural ingredient not present in bare general relativity: a choice of clock subsystem, a choice of statistical state, a choice of quantum frame, or an explicit new postulated parameter, respectively. Practical global timekeeping sidesteps the entire foundational debate by fixing one such extra ingredient (the geoid, or the barycentric frame) by international convention and simply propagating relativistic corrections through it with extraordinary precision.

%======================================================================
\section{Open Problems and Outlook}
\label{sec:outlook}
%======================================================================

Several fronts remain genuinely unsettled as of 2026:

\begin{itemize}
\item \textbf{Uniqueness and robustness of relational time.} The equivalence results unifying Page--Wootters, Dirac quantization, and deparametrization (the ``trinity of relational dynamics'') hold under specific regularity assumptions (e.g., idealized, often monotonic, clocks); extending them to realistic, non-ideal, interacting, or periodic clocks, and to full non-perturbative quantum gravity (as opposed to minisuperspace or group-field-theory toy models), is ongoing work.
\item \textbf{Experimental access to quantum time.} Proposals to witness a genuine superposition of proper times, or to measure the thermodynamic cost of extracting classical time from a quantum clock, are now technically within reach of atom-interferometric and trapped-ion platforms; whether near-term experiments can cleanly distinguish these effects from more mundane decoherence is not yet settled.
\item \textbf{Thermal time beyond cosmology.} The single clean success of the thermal time hypothesis (recovering FLRW cosmic time from the CMB state) has not been matched by an equally clean account of thermal time far from equilibrium, or by a fully worked-out, matter-\emph{and}-gravity version of the construction.
\item \textbf{The status of external parameters such as SHP's $\tau$.} Parametrized theories with a strictly external, unobservable universal time raise a foundational question that is the mirror image of the problem of time: if $\tau$ is in principle unobservable, what distinguishes the claim that it is physically real from the claim that it is a convenient but dispensable calculational device, and can this be settled empirically (for instance, via off-shell or pair-production phenomenology) rather than only interpretively?
\item \textbf{Whether ``global time'' is the right question.} A recurring theme across quantum-reference-frame and relational-dynamics work is that many puzzles dissolve, rather than resolve, once one stops asking for a single global time (or frame) altogether and instead asks only for well-defined \emph{relative} quantities between subsystems; whether this perspectival, relational stance is the correct \emph{final} lesson of the subject, or an intermediate way-station on the way to a still-undiscovered global structure, remains the deepest open question surveyed in this review.
\end{itemize}

%======================================================================
\section*{Acknowledgment of Scope}
%======================================================================
\addcontentsline{toc}{section}{Acknowledgment of Scope}
This review is necessarily selective. It omits, for reasons of space, loop quantum cosmology's specific treatment of the big-bang singularity and internal time, causal set theory's discrete account of temporal becoming, Barbour-style shape dynamics and Machian relationalism (in which refoliation invariance is traded for local conformal invariance and time is eliminated from the configuration space itself), and the extensive philosophy-of-physics literature on the A-theory/B-theory of time and presentism/eternalism as they interact with relativity. Each is a legitimate and active part of the modern landscape and a natural next step for a longer treatment.

%======================================================================
\begin{thebibliography}{99}
%======================================================================

\bibitem{Malament1977} D.~B. Malament, ``Causal Theories of Time and the Conventionality of Simultaneity,'' \textit{No\^us} \textbf{11}, 293 (1977).

\bibitem{Geroch1970} R. Geroch, ``Domain of Dependence,'' \textit{J. Math. Phys.} \textbf{11}, 437 (1970).

\bibitem{BernalSanchez2003} A.~N. Bernal and M. S\'anchez, ``On Smooth Cauchy Hypersurfaces and Geroch's Splitting Theorem,'' \textit{Commun. Math. Phys.} \textbf{243}, 461 (2003).

\bibitem{BernalSanchez2005} A.~N. Bernal and M. S\'anchez, ``Smoothness of Time Functions and the Metric Splitting of Globally Hyperbolic Spacetimes,'' \textit{Commun. Math. Phys.} \textbf{257}, 43 (2005).

\bibitem{BernalSanchez2007} A.~N. Bernal and M. S\'anchez, ``Globally Hyperbolic Spacetimes Can Be Defined as `Causal' Instead of `Strongly Causal','' \textit{Class. Quantum Grav.} \textbf{24}, 745 (2007).

\bibitem{AGH1998} L. Andersson, G.~J. Galloway, and R. Howard, ``The Cosmological Time Function,'' \textit{Class. Quantum Grav.} \textbf{15}, 309 (1998).

\bibitem{Sanchez2022} M. S\'anchez, ``Globally Hyperbolic Spacetimes: Slicings, Boundaries and Counterexamples,'' \textit{Gen. Relativ. Gravit.} \textbf{54}, 124 (2022).

\bibitem{Kuchar1992} K.~V. Kucha\v{r}, ``Time and Interpretations of Quantum Gravity,'' in \textit{Proc. 4th Canadian Conf. on General Relativity and Relativistic Astrophysics} (World Scientific, Singapore, 1992).

\bibitem{Isham1993} C.~J. Isham, ``Canonical Quantum Gravity and the Problem of Time,'' in \textit{Integrable Systems, Quantum Groups, and Quantum Field Theories}, arXiv:gr-qc/9210011 (1993).

\bibitem{Anderson2012} E. Anderson, ``The Problem of Time in Quantum Gravity,'' arXiv:1009.2157 [gr-qc].

\bibitem{ClockField2026} A.~G. Fernandes de Freitas \textit{et al.}, ``Geometric Emergence of Time in Canonical Quantum Gravity,'' \textit{Class. Quantum Grav.} \textbf{43}, 115004 (2026).

\bibitem{Reassessing2023} (author redacted in source listing), ``Reassessing the Problem of Time of Quantum Gravity,'' \textit{Gen. Relativ. Gravit.} (2023).

\bibitem{GRF2025} Essay submitted to the Gravity Research Foundation 2025 Awards, ``The Problem of Time and its Quantum Resolution,'' arXiv:2504.00152.

\bibitem{PageWootters1983} D.~N. Page and W.~K. Wootters, ``Evolution Without Evolution: Dynamics Described by Stationary Observables,'' \textit{Phys. Rev. D} \textbf{27}, 2885 (1983).

\bibitem{HohnVanrietvelde2021} P.~A. H\"ohn, A.~R.~H. Smith, and M.~P.~E. Lock, ``Trinity of Relational Quantum Dynamics,'' \textit{Phys. Rev. D} \textbf{104}, 066001 (2021).

\bibitem{ChataignierEtAl2024} L. Chataignier, P.~A. H\"ohn, M.~P.~E. Lock, and F.~M. Mele, ``Relational Dynamics with Periodic Clocks,'' \textit{New J. Phys.} \textbf{28}, 034504 (2026); arXiv:2409.06479.

\bibitem{CalcinariGielen2025} A. Calcinari and S. Gielen, ``Relational Dynamics and Page--Wootters Formalism in Group Field Theory,'' \textit{Quantum} \textbf{9}, 1610 (2025).

\bibitem{AresEtAl2025} T. \textit{et al.} (P. Erker, N. Ares and collaborators), ``Entropic Costs of Extracting Classical Ticks from a Quantum Clock,'' \textit{Phys. Rev. Lett.} \textbf{135}, 200407 (2025).

\bibitem{FavalliEtAl2025} T. Favalli \textit{et al.}, on the emergence of time and space from entanglement in an extended Page--Wootters model, Univ.~of Naples ``Federico~II'' preprint (2025).

\bibitem{ConnesRovelli1994} A. Connes and C. Rovelli, ``Von Neumann Algebra Automorphisms and Time-Thermodynamics Relation in Generally Covariant Quantum Theories,'' \textit{Class. Quantum Grav.} \textbf{11}, 2899 (1994).

\bibitem{Rovelli1993} C. Rovelli, ``Statistical Mechanics of Gravity and the Thermodynamical Origin of Time,'' \textit{Class. Quantum Grav.} \textbf{10}, 1549 (1993).

\bibitem{ThermalTimePOVM2024} (authors as listed in source), ``Thermal Time as an Unsharp Observable,'' \textit{J. Math. Phys.} \textbf{65}, 032105 (2024).

\bibitem{Swanson2021} N. Swanson, ``Can Quantum Thermodynamics Save Time?,'' \textit{Philos. Sci.} \textbf{88}, 205 (2021).

\bibitem{GiacominiCastroRuizBrukner2019} F. Giacomini, E. Castro-Ruiz, and \v{C}. Brukner, ``Quantum Mechanics and the Covariance of Physical Laws in Quantum Reference Frames,'' \textit{Nat. Commun.} \textbf{10}, 494 (2019).

\bibitem{Giacomini2021} F. Giacomini, ``Spacetime Quantum Reference Frames and Superpositions of Proper Times,'' \textit{Quantum} \textbf{5}, 508 (2021).

\bibitem{SmithAhmed2020} A.~R.~H. Smith and M. Ahmadi, ``Quantum Clocks Observe Classical and Quantum Time Dilation,'' \textit{Nat. Commun.} \textbf{11}, 5360 (2020).

\bibitem{deLaHametteEtAl2025} A.-C. de la Hamette, V. Kabel, M. Christodoulou, and \v{C}. Brukner, ``Indefinite Causal Order and Quantum Coordinates,'' \textit{Phys. Rev. Lett.} \textbf{135}, 141402 (2025).

\bibitem{KabelEtAl2025} V. Kabel, A.-C. de la Hamette, L. Apadula, C. Cepollaro, H. Gomes, J. Butterfield, and \v{C}. Brukner, ``Quantum Coordinates, Localisation of Events, and the Quantum Hole Argument,'' \textit{Commun. Phys.} \textbf{8}, 185 (2025).

\bibitem{SorciEtAl2026} G. Sorci, J. Foo, D. Leibfried, C. Sanner, and I. Pikovski, ``Quantum Signatures of Proper Time in Optical Ion Clocks,'' \textit{Phys. Rev. Lett.} \textbf{136}, 163602 (2026).

\bibitem{Stueckelberg1941} E.~C.~G. Stueckelberg, ``La Signification du Temps Propre en M\'ecanique Ondulatoire,'' \textit{Helv. Phys. Acta} \textbf{14}, 322 (1941).

\bibitem{Stueckelberg1942} E.~C.~G. Stueckelberg, ``La M\'ecanique du Point Mat\'eriel en Th\'eorie de Relativit\'e et en Th\'eorie des Quanta,'' \textit{Helv. Phys. Acta} \textbf{15}, 23 (1942).

\bibitem{HorwitzPiron1973} L.~P. Horwitz and C. Piron, ``Relativistic Dynamics,'' \textit{Helv. Phys. Acta} \textbf{46}, 316 (1973).

\bibitem{Fanchi1993} J.~R. Fanchi, ``Review of Invariant Time Formulations of Relativistic Quantum Theories,'' \textit{Found. Phys.} \textbf{23}, 487 (1993).

\bibitem{Horwitz2015} L.~P. Horwitz, \textit{Relativistic Quantum Mechanics} (Springer, Dordrecht, 2015).

\bibitem{Horwitz2019} L.~P. Horwitz, ``An Elementary Canonical Classical and Quantum Dynamics for General Relativity,'' \textit{Eur. Phys. J. Plus} \textbf{134}, 213 (2019).

\bibitem{Land2020} M. Land, ``A 4+1 Formalism for the Evolving Stueckelberg--Horwitz--Piron Metric,'' \textit{Symmetry} \textbf{12}, 1721 (2020).

\bibitem{HorwitzArshanskyElitzur1988} L.~P. Horwitz, R.~I. Arshansky, and A.~C. Elitzur, ``On the Two Aspects of Time: The Distinction and Its Implications,'' \textit{Found. Phys.} \textbf{18}, 1159 (1988).

\bibitem{CosmicTimeQM2022} (authors as listed in source), ``Cosmic-Time Quantum Mechanics and the Passage-of-Time Problem,'' arXiv:2209.06070.

\bibitem{CGPM2022} 27th General Conference on Weights and Measures (CGPM), Resolution 5, ``On the Future Redefinition of the Second'' (2022), \url{https://www.bipm.org/en/cgpm-2022/resolution-5}.

\bibitem{IAU2000B19} International Astronomical Union, Resolution B1.9 (2000), on redefinition of Terrestrial Time; and related IAU 1991/2006 resolutions defining TCG, TCB, and TT.

\bibitem{TavellaPetit2020} P. Tavella and G. Petit, ``Precise Time Scales and Navigation Systems: Mutual Benefits of Timekeeping and Positioning,'' \textit{Satell. Navig.} \textbf{1}, 1 (2020).

\bibitem{Takano2016} T. Takano, M. Takamoto, I. Ushijima \textit{et al.}, ``Geopotential Measurements with Synchronously Linked Optical Lattice Clocks,'' \textit{Nat. Photonics} \textbf{10}, 662 (2016).

\end{thebibliography}

\end{document}
